%% ma: L12-B15-0013
%% lop: 12
%% chuong: 5
%% bai: 15
%% dang: 12.15.11
%% loai: TN4
%% muc: NB
%% dap_an: "A"
%% nguon: "(NBV)-ĐỀ SỐ 16-MỖI NGÀY 1 ĐỀ THI 2025.pdf (D16, MỖI NGÀY 1 ĐỀ - Đề số 16), Phần I Câu 12"
%% kien_thuc: "Bài 14 (vectơ pháp tuyến của mặt phẳng), Bài 15 (vectơ chỉ phương của đường thẳng)"
%% trang_thai: cho-duyet
%% ly_do: "Dạng mới đề xuất 12.15.11 chờ giáo viên duyệt"
%% ghi_chu: "Đáp án gốc A đúng."
\begin{ex}
Trong không gian $Oxyz$, cho đường thẳng $d:\dfrac{x-1}{2}=\dfrac{y-2}{-1}=\dfrac{z-3}{2}$. Mặt phẳng $(P)$ vuông góc với đường thẳng $d$ có một vectơ pháp tuyến là
\choice{\True $\vec n_2=(2;-1;2)$}{$\vec n_3=(-1;-2;-3)$}{$\vec n_1=(1;2;3)$}{$\vec n_4=(2;1;2)$}
\loigiai{Đường thẳng $d$ có vectơ chỉ phương $\vec u=(2;-1;2)$. Vì $(P)\perp d$ nên $\vec u$ cũng là một vectơ pháp tuyến của $(P)$. Chọn A.}
\end{ex}
